\documentclass[10pt,a4paper]{amsart}
\usepackage[margin=19mm]{geometry}
\usepackage{amsmath,amssymb,mathtools}
\usepackage[scr=rsfs,cal=euler]{mathalfa}
\usepackage{xfrac}
\usepackage[unicode,hidelinks,bookmarks=false]{hyperref}
\newcommand{\ie}{i.e.\ }
\newcommand{\cf}{cf.\ }
\newcommand{\resp}{respectively\ }
\newcommand{\Subsection}{\subsection}
\providecommand{\nobreakdash}{\nobreak\hbox{-}}
\setlength{\emergencystretch}{2em}
\multlinegap0pt
\newcommand{\Split}{\mathrm{Split}}
\newcommand{\inte}{\mathrm{int}\,}
\newcommand\labs[1]{\left\lvert#1\right\rvert}
\theoremstyle{plain}
\newtheorem{thm}{Theorem}[section]
\newtheorem{lem}[thm]{Lemma}
\newtheorem{pro}[thm]{Proposition}
\newtheorem{cor}[thm]{Corollary}
\newtheorem{claim}[thm]{Claim}
\newtheorem{qu}[thm]{Question}
\newtheorem{prb}{Problem}
\newtheorem{con}{Conjecture}
\newtheorem{thmAlph}{Theorem}[section]
\newtheorem{corAlph}[thmAlph]{Corollary}
\theoremstyle{definition}
\newtheorem{exam}[thm]{Example}
\newtheorem{defn}[thm]{Definition}
\theoremstyle{remark}
\newtheorem{rema}[thm]{Remark}
\newtheorem*{ack}{Acknowledgements}
\title[Ramsey Theory for Product Trees]{Ramsey Theory for Product Trees}
\author{Alexander Fish, Sean Skinner}
\date{Source: arXiv:2609.28898v1 (CC BY 4.0)}
\begin{document}
\maketitle

\noindent\emph{Source and licence.} The delivered work is the article shown in the title above, version arXiv:2609.28898v1 (CC BY 4.0), distributed under the Creative Commons Attribution 4.0 International licence, as recorded in the arXiv licence metadata for this exact version and shown on the abstract page saved with this item. The full licence notice, which carries the licence URL and the address of that abstract page, is the bundled file \texttt{licenses/arxiv-2609.28898-CC-BY-4.0.txt}.
\par\smallskip

\noindent\textit{Editorial scope.} Counted unit: the article's Pro (source label Prop: 2d binary splitting count) (source0.tex lines 491--494, source label \texttt{Prop: 2d binary splitting count}): ``For all and we have that 2 (2) (A) A -(2 n 2 m -1).''. It is delivered together with the article's complete original proof (source0.tex lines 506--527, one proof environment printed later in the article, detached from the statement by the article's own arrangement, 22 lines), copied line for line from the article's own text. Required context retained uncounted: Lemma: Baby split lem (lines 496--500). Editorial formatting only: an AMS \texttt{amsart} preamble that restates the article's own macro definitions and its own theorem declarations, and the theorem environments the retained lines use; the article's own class file is neither shipped nor input; the retained environments are renumbered sequentially in order of appearance, whereas the article prints them with its own numbering; the article's equation numbering is not reproduced and every cross-reference inside the retained text resolves to the wrapper's numbering. No citation occurs inside any retained line, so the wrapper carries no bibliography; All retained bodies are exact copies of the bundled \texttt{sources/211625/source0.tex}, so no omission receipt applies. No asset-inclusion command occurs inside any retained span and the package ships no image asset.


\section*{Retained context: Lemma: Baby split lem (source0.tex lines 496--500)}

\begin{lem}\label{Lemma: Baby split lem}
For all $n\geq 0$ and $C\subset \partial T_n$ we have that
\[ \labs{\Split_1^{(2)}(C)} = \labs{C} - \mathbf{1}_{C\neq \emptyset}.\]
In particular, if $C\neq \emptyset$, then $T_C$ has exactly $\labs{C}-1$ vertices with $2$ children.
\end{lem}


\section{The counted unit: Pro (source label Prop: 2d binary splitting count) and its complete original proof}

\begin{pro}\label{Prop: 2d binary splitting count}
For all $n,m\geq 0$ and $A \subset \partial \Gamma_{n,m}$ we have that
\[ \labs{\Split_2^{(2)}(A)} \geq \labs{A} -(2^n + 2^m -1).\] 
\end{pro}

\begin{proof}[Proof of Proposition~\ref{Prop: 2d binary splitting count}]
Let $A\subset \partial \Gamma_{n,m}$. For each $x \in \{0,1\}^n$ let
\[ A_x:= \{y \in \{0,1\}^m \, : \, (x,y) \in A\}\]
be the vertical fibre above $x$. For each $v \in \inte T_m$ consider
\[ E_v:= \{ x \in \{0,1\}^n \, : \, v \in \Split_1^{(2)}(A_x)\} \subset \partial T_n\]
as a one-dimensional leaf set. By counting in two orders and using Lemma~\ref{Lemma: Baby split lem} we have that
\begin{align} 
\sum_{v\in \inte T_m} \labs{E_v} &= \sum_{x\in \{0,1\}^n} \labs{\Split_1^{(2)}(A_x)} \nonumber \\
&= \sum_{x\in \{0,1\}^n} \left( \labs{A_x} - \mathbf{1}_{A_x \neq \emptyset}\right)\label{eq: first use of baby split}
\geq \labs{A} - 2^n.
\end{align}
Let $u \in \Split_1^{(2)}(E_v)$ for some $v\in \inte T_m$. We claim that $(u,v) \in \Split_2^{(2)}(A)$.
Indeed, by definition of $\Split_1^{(2)}(E_v)$, for each $\varepsilon \in \{0,1\}$ there exists some $x_\varepsilon \in E_v$ such that $u\varepsilon \preceq x_\varepsilon$. By definition of $E_v$, for each $\varepsilon \in \{0,1\}$, $v \in \Split_1^{(2)}(A_{x_\varepsilon})$, and so for each $\delta \in \{0,1\}$ there exists some $y_{\delta}^{\varepsilon} \in A_{x_\varepsilon}$ such that $v\delta \preceq y_{\delta}^{\varepsilon}$. It then follows that $(u\varepsilon,v\delta) \preceq_2 (x_\varepsilon,y_{\delta}^\varepsilon) \in A$ for each $(\varepsilon,\delta )\in \{0,1\}^2$, which proves the claim.

It follows that
\begin{align*}
\labs{\Split_2^{(2)}(A)} &\geq \sum_{v\in  \inte T_m} \labs{\Split_1^{(2)}(E_v)}\\
&\geq\sum_{v\in \inte T_m} \left(\labs{E_v}-1\right)\\
&\geq \labs{A}-2^n - (1+2+\cdots+2^{m-1}) = \labs{A} - (2^n + 2^m-1).
\end{align*}
where in the second line we use Lemma~\ref{Lemma: Baby split lem} again, noting that the bound $\labs{\Split_1^{(2)}(E_v) }\geq \labs{E_v}-1$ holds irrespective of whether $E_v$ is empty or not, and in the last line we use equation~\eqref{eq: first use of baby split}.
\end{proof}

\end{document}
